% !TEX program = pdflatex
\documentclass[11pt,a4paper]{article}

% =========================================================================
% GUÍA DE SUSTENTACIÓN DE LA UNIDAD I
% Proyecto: CasaScore (AlquilaPuno)
% Gestión de Proyectos TICS (EST315) -- Semestre 2026-II -- Team Sobrevivientes
% =========================================================================

% --- PAQUETES DE CONFIGURACIÓN ---
\usepackage[utf8]{inputenc}
\usepackage[T1]{fontenc}
\usepackage[spanish,es-tabla]{babel}
% Compatibilidad con babel-spanish reciente (flechas de TikZ y signo de porcentaje)
\AtBeginDocument{\shorthandoff{<>}\spanishplainpercent}
\usepackage{geometry}
\geometry{
    top=3.6cm,
    bottom=2.3cm,
    left=2.2cm,
    right=2.2cm,
    % PARTE MODIFICABLE: headheight debe ser mayor que la altura del logo más alto (2.56cm = 73pt).
    headheight=80pt,
    headsep=15pt,
    footskip=20pt
}

\usepackage{graphicx}
\usepackage{fancyhdr}
\usepackage[table]{xcolor}
\usepackage{titlesec}
\usepackage{booktabs}
\usepackage{tabularx}
\usepackage{array}
\usepackage{amsmath}
\usepackage{enumitem}
\usepackage{microtype}
\usepackage{amssymb}
% --- RECURSOS VISUALES: GRÁFICOS Y CÓDIGO DE COLORES ---
\usepackage{tikz}
\usetikzlibrary{positioning,calc}
\definecolor{GrisFila}{RGB}{246, 248, 251}
\definecolor{CuadCerca}{RGB}{150, 30, 30}
\definecolor{CuadSatisf}{RGB}{180, 130, 20}
\definecolor{CuadInfo}{RGB}{10, 45, 90}
\definecolor{CuadMonit}{RGB}{110, 115, 125}
\newcommand{\etiqueta}[2]{{\setlength{\fboxsep}{2pt}\colorbox{#1}{\textcolor{white}{\scriptsize\bfseries #2}}}}
\newcommand{\cerca}{\etiqueta{CuadCerca}{Gestionar de cerca}}
\newcommand{\satisf}{\etiqueta{CuadSatisf}{Mantener satisfecho}}
\newcommand{\informado}{\etiqueta{CuadInfo}{Mantener informado}}
\newcommand{\monit}{\etiqueta{CuadMonit}{Monitorear}}
\usepackage{needspace}
\usepackage{longtable}
\definecolor{Rama1}{RGB}{10, 45, 90}
\definecolor{Rama2}{RGB}{40, 95, 150}
\definecolor{Rama3}{RGB}{180, 130, 20}
\definecolor{Rama4}{RGB}{30, 110, 60}
\definecolor{Rama5}{RGB}{110, 60, 140}
\definecolor{RaciR}{RGB}{150, 30, 30}
\definecolor{RaciA}{RGB}{10, 45, 90}
\definecolor{RaciC}{RGB}{180, 130, 20}
\definecolor{RaciI}{RGB}{140, 145, 155}
\definecolor{NivAlto}{RGB}{30, 110, 60}
\definecolor{NivMedio}{RGB}{180, 130, 20}
\definecolor{NivBajo}{RGB}{150, 30, 30}
\usetikzlibrary{patterns}
\definecolor{SevBaja}{RGB}{30, 110, 60}
\definecolor{SevMedia}{RGB}{180, 130, 20}
\definecolor{SevAlta}{RGB}{150, 30, 30}
\usetikzlibrary{arrows.meta,shapes.geometric,fit,backgrounds}
\usepackage{listings}
\lstset{basicstyle=\ttfamily\scriptsize, frame=single, rulecolor=\color{AzulUNA!40}, backgroundcolor=\color{GrisFila}, columns=fullflexible, keepspaces=true, breaklines=true, xleftmargin=4pt, xrightmargin=4pt, literate={á}{{\'a}}1 {é}{{\'e}}1 {í}{{\'i}}1 {ó}{{\'o}}1 {ú}{{\'u}}1 {ñ}{{\~n}}1 {Á}{{\'A}}1}

\usepackage{hyperref}

% Logos: funciona a 2 o 3 niveles de profundidad dentro de docs/latex/
\graphicspath{{../../assets/img/}{../../../assets/img/}}

% --- PALETA INSTITUCIONAL FIEI / UNA PUNO ---
\definecolor{AzulUNA}{RGB}{10, 45, 90}
\definecolor{DoradoFIEI}{RGB}{180, 130, 20}
\definecolor{GrisOscuro}{RGB}{50, 50, 50}
\definecolor{AzulClaro}{RGB}{228, 236, 246}

% --- CONFIGURACIÓN DE LISTAS ---
\setlist{topsep=2pt, itemsep=2pt, parsep=1pt}

% --- COLUMNAS DE TABLA ---
\newcolumntype{L}{>{\raggedright\arraybackslash}X}
\newcolumntype{P}[1]{>{\raggedright\arraybackslash}p{#1}}
\newcolumntype{C}[1]{>{\centering\arraybackslash}p{#1}}
\newcolumntype{Q}[1]{>{\raggedleft\arraybackslash}p{#1}}
\newcolumntype{R}{>{\raggedleft\arraybackslash}X}
\renewcommand{\arraystretch}{1.25}

% --- FORMATO DE ENLACES ---
\hypersetup{
    colorlinks=true,
    linkcolor=AzulUNA,
    urlcolor=AzulUNA,
    citecolor=AzulUNA,
    pdftitle={Guía de Sustentación de la Unidad I -- CasaScore (AlquilaPuno)},
    pdfauthor={Team Sobrevivientes}
}

% --- FORMATO DE TÍTULOS ---
\titleformat{\section}{\color{AzulUNA}\normalfont\large\bfseries}{\thesection.}{0.7em}{}[\color{AzulUNA}\titlerule]
\titleformat{\subsection}{\color{AzulUNA}\normalfont\normalsize\bfseries}{\thesubsection.}{0.5em}{}
\titlespacing*{\section}{0pt}{9pt}{4pt}
\titlespacing*{\subsection}{0pt}{6pt}{2pt}

% --- ENCABEZADO Y PIE DE PÁGINA (fancyhdr) ---
\pagestyle{fancy}
\fancyhf{}

% ===== TAMAÑO DE LOS LOGOS (PARTE MODIFICABLE) =====
% logo_fiei.png  : height=1.45cm (lado izquierdo)
% logo_equipo.png: height=2.56cm (lado derecho; 25 % más grande porque la imagen tiene más borde blanco)
% Para cambiar el tamaño, modifique solo el valor height=... de cada \includegraphics.
% Si el logo más alto supera 2.6cm, aumente también headheight y top en la configuración de página.
\fancyhead[L]{%
    \includegraphics[height=1.45cm,keepaspectratio]{logo_fiei.png}%
}
\fancyhead[C]{%
    \scriptsize \textbf{\textcolor{AzulUNA}{UNIVERSIDAD NACIONAL DEL ALTIPLANO DE PUNO}}\\[2pt]
    \tiny \textcolor{GrisOscuro}{FACULTAD DE INGENIERÍA ESTADÍSTICA E INFORMÁTICA}\\[2pt]
    \scriptsize \textbf{\textcolor{AzulUNA}{Gestión de Proyectos TICS (EST315) \textbar\ Team Sobrevivientes}}\\[2pt]
    \footnotesize \textbf{\textcolor{DoradoFIEI}{GUÍA DE SUSTENTACIÓN DE LA UNIDAD I}}%
}
\fancyhead[R]{%
    \includegraphics[height=2.56cm,keepaspectratio]{logo_equipo.png}%
}

\renewcommand{\headrulewidth}{0pt}
\renewcommand{\headrule}{}

\fancyfoot[L]{\scriptsize \textcolor{GrisOscuro}{CasaScore (AlquilaPuno)}}
\fancyfoot[C]{\small \thepage}
\fancyfoot[R]{\scriptsize \textcolor{GrisOscuro}{Puno, Perú -- Semestre 2026-II}}
\renewcommand{\footrulewidth}{0.4pt}

\begin{document}

\begin{center}
    {\Large \textbf{\textcolor{AzulUNA}{GUÍA DE SUSTENTACIÓN DE LA UNIDAD I}}}\\[3pt]
    {\normalsize \textbf{CasaScore (AlquilaPuno)}}\\[4pt]
    \footnotesize \textbf{Fecha:} semana 8 (12/10 -- 16/10/2026) \quad \textbar\quad \textbf{Versión:} 1.0 \quad \textbar\quad \textbf{Fase:} Implementación
\end{center}

\vspace{0.1cm}

\section{Datos Generales}
\noindent\begin{tabularx}{\textwidth}{@{}P{4.2cm} L@{}}
\toprule
\textbf{Proyecto:} & CasaScore (aplicación AlquilaPuno) \\
\textbf{Responsable:} & Clyde Neil Paricahua Pari -- Project Manager \\
\textbf{Proceso PM4DEV:} & Gestión de la información (comunicación de resultados) \\
\textbf{Modalidad:} & Exposición oral con diapositivas y demostración de la app, seguida de preguntas \\
\textbf{Duración prevista:} & 15 minutos de exposición y 5 minutos de preguntas \\
\textbf{Expositores:} & Los cuatro integrantes del equipo \\
\bottomrule
\end{tabularx}

\section{Mensaje Central}
\begin{center}
\begin{minipage}{0.88\textwidth}\centering\itshape
CasaScore ayuda a cualquier persona que busca cuarto en Puno a decidir con un puntaje de 0 a 100 que combina precio, cercanía, seguridad, servicios y tipo de baño según sus propias prioridades, y que explica cada punto sumado o restado.
\end{minipage}
\end{center}

\Needspace{16\baselineskip}
\section{Estructura de la Presentación}
{\footnotesize
\setlength{\tabcolsep}{3.5pt}
\rowcolors{2}{white}{GrisFila}
\noindent\begin{tabularx}{\textwidth}{@{}C{0.6cm} P{2.9cm} L P{2.7cm} C{1.1cm}@{}}
\toprule
\rowcolor{AzulClaro}
\textbf{N.\textsuperscript{o}} & \textbf{Diapositiva} & \textbf{Mensaje principal} & \textbf{Expone} & \textbf{Min.} \\
\midrule
1 & Portada & Proyecto, equipo, curso y asesor. & Clyde N. Paricahua Pari & 0.5 \\
2 & El problema & Buscar cuarto en Puno: anuncios dispersos, datos sin verificar y reglas de convivencia ocultas. & Clyde N. Paricahua Pari & 1.0 \\
3 & La solución & CASA SCORE: puntaje de 0 a 100 explicado, según las prioridades y el punto de referencia de cada arrendatario. & Clyde N. Paricahua Pari & 1.0 \\
4 & Cartera y decisión & Dos proyectos evaluados; Horus Sentinel excluido en la semana 5 porque necesitaba 19 semanas y S/ 13,663.57. & Clyde N. Paricahua Pari & 1.0 \\
5 & Interesados & Matriz poder--interés: arrendatarios y propietarios se gestionan de cerca. & Viviana M. Quisbert Quispe & 0.5 \\
6 & Alcance y EDT & 5 ramas, 23 paquetes y 512 horas (regla del 100\,\%). & Ruth K. Apaza Solis & 1.0 \\
7 & Cronograma & Ruta crítica, PERT y fin previsto el 16/12/2026 con 75\,\% de probabilidad de cumplir la semana 17. & Ruth K. Apaza Solis & 1.5 \\
8 & Presupuesto y viabilidad & S/ 12,829.82 en total; S/ 626.86 de aporte por integrante; VAN económico positivo. & Ruth K. Apaza Solis & 1.5 \\
9 & Riesgos y calidad & 3 riesgos altos con respuesta; ISO/IEC 25010 y 14 casos de prueba. & Viviana M. Quisbert Quispe & 1.5 \\
10 & Arquitectura & App Android en tres capas, motor en el celular y Supabase como backend. & Eliseo Tarqui Ajahuana & 1.0 \\
11 & Demostración & Ranking, cambio de perfil, detalle del puntaje y publicación de un cuarto. & Eliseo Tarqui Ajahuana & 2.5 \\
12 & Estado y control & SPI y CPI al corte, solicitud de cambio rechazada y plan de recuperación. & Ruth K. Apaza Solis & 1.0 \\
13 & Próximos pasos & Supabase, punto de referencia, piloto en la semana 13 y MVP final en la semana 16. & Clyde N. Paricahua Pari & 1.0 \\
\midrule
\rowcolor{white}\multicolumn{4}{@{}l}{\textbf{Tiempo total de exposición}} & \textbf{15} \\
\bottomrule
\end{tabularx}
}

\vspace{0.1cm}
\noindent\textbf{Tiempo de cada integrante}
\begin{center}
\begin{tikzpicture}[font=\scriptsize]
\fill[AzulUNA] (0.00,0) rectangle (4.20,0.55); \node[text=white, font=\tiny\bfseries, align=center] at (2.10,0.275) {Clyde\\ 4.5 min};
\fill[Rama4] (4.20,0) rectangle (6.07,0.55); \node[text=white, font=\tiny\bfseries, align=center] at (5.13,0.275) {Viviana\\ 2.0 min};
\fill[Rama2] (6.07,0) rectangle (10.73,0.55); \node[text=white, font=\tiny\bfseries, align=center] at (8.40,0.275) {Ruth\\ 5.0 min};
\fill[Rama3] (10.73,0) rectangle (14.00,0.55); \node[text=white, font=\tiny\bfseries, align=center] at (12.37,0.275) {Eliseo\\ 3.5 min};
\end{tikzpicture}
\end{center}

\Needspace{16\baselineskip}
\section{Cifras Clave}
{\small
\rowcolors{2}{white}{GrisFila}
\noindent\begin{tabularx}{\textwidth}{@{}P{4.2cm} L P{3.2cm}@{}}
\toprule
\rowcolor{AzulClaro}
\textbf{Dato} & \textbf{Valor} & \textbf{Documento de origen} \\
\midrule
Esfuerzo del proyecto & 512 horas del equipo (8 h semanales por integrante, 16 semanas) & EDT \\
Presupuesto total & S/ 12,829.82 (línea base S/ 11,663.47 y reserva de gestión S/ 1,166.35) & Presupuesto \\
Gasto en efectivo & S/ 2,507.44; S/ 626.86 por integrante & Presupuesto \\
Evaluación a 3 años & VAN financiero negativo (S/ $-$17,077.02); VAN económico positivo (S/ 64,059.98) & Viabilidad \\
Fin previsto & 16/12/2026; 75\,\% de probabilidad de terminar hasta el 18/12 & Cronograma \\
Desempeño al corte (semana 6) & SPI 0.56 y CPI 0.96: dos semanas de atraso con costo controlado & Control del proyecto \\
Riesgos altos & Pocos propietarios, retraso en Supabase y baja adopción en el piloto & Plan de riesgos \\
Calidad & 7 casos de prueba del motor aprobados; meta de satisfacción de 80\,\% & Validación del producto \\
Horus Sentinel & 19 semanas requeridas y S/ 13,663.57; excluido el 22/09/2026 & Acta de exclusión \\
\bottomrule
\end{tabularx}
}

\Needspace{14\baselineskip}
\section{Guion de la Demostración}
\begin{enumerate}[leftmargin=*]
    \item Abrir la app con los 5 cuartos de prueba: el ranking aparece ordenado de mayor a menor puntaje.
    \item Cambiar del perfil por defecto al perfil de cercanía: el segundo lugar pasa de Tesista a Laykakota.
    \item Abrir el detalle del cuarto Bellavista: mostrar que tiene el mejor precio, pero baja a 8 por sus penalizaciones.
    \item Abrir el detalle del cuarto de Puerta 1: mostrar sus bonificaciones y su puntaje de 100.
    \item Publicar un cuarto nuevo con sus 4 fotos y verlo ubicado en el ranking.
\end{enumerate}
\noindent\textbf{Plan alternativo:} si la app o la conexión fallan, se reproduce un video de 2 minutos grabado el día anterior con los mismos pasos.

\Needspace{16\baselineskip}
\section{Preguntas Probables}
{\footnotesize
\setlength{\tabcolsep}{3.5pt}
\rowcolors{2}{white}{GrisFila}
\noindent\begin{tabularx}{\textwidth}{@{}P{4.0cm} L P{2.5cm}@{}}
\toprule
\rowcolor{AzulClaro}
\textbf{Pregunta probable} & \textbf{Respuesta} & \textbf{Responde} \\
\midrule
Por qué un puntaje y no solo filtros? & Los filtros descartan cuartos; el puntaje los compara en una sola escala de 0 a 100 y muestra por qué cada cuarto tiene ese valor. & Eliseo Tarqui Ajahuana \\
Cómo evitan que un propietario mienta? & Fotos obligatorias de fachada, habitación, baño y ventana, y una penalización de 25 puntos si las condiciones no estaban declaradas. En el año 1 se agregan reportes y calificaciones con estrellas. & Eliseo Tarqui Ajahuana \\
Por qué excluyeron Horus Sentinel? & Su cronograma requería 19 semanas y solo quedaban 13; además, sumado a CasaScore, cada integrante debía aportar S/ 1,434.19. Se decidió por unanimidad el 22/09. & Clyde N. Paricahua Pari \\
El proyecto es rentable? & En lo financiero no en los 3 primeros años; en lo económico sí, por el ahorro de tiempo y dinero de los arrendatarios (VAN económico positivo). Por eso se busca financiamiento en concursos. & Ruth K. Apaza Solis \\
Por qué van atrasados y qué hacen? & La gestión acumuló dos semanas de atraso por la carga de dos proyectos. Se rechazó aumentar horas y se aplica un plan de recuperación sin costo adicional. & Ruth K. Apaza Solis \\
Cómo saben que el ranking es correcto? & El motor está separado de las pantallas y se probó con 7 casos y 5 cuartos de prueba; el orden obtenido coincide con el esperado. & Viviana M. Quisbert Quispe \\
Qué pasa con los datos personales? & Solo se guardan los datos necesarios, se pide consentimiento al registrarse y cada usuario solo ve lo que le corresponde (Ley N.\textsuperscript{o} 29733). & Viviana M. Quisbert Quispe \\
Cuánto cuesta operar la app? & El MVP funciona con planes gratuitos; el único pago es la cuenta de Google Play (S/ 85.90). & Eliseo Tarqui Ajahuana \\
Qué aplicaron del PMBOK y de PM4DEV? & Las 6 fases y los 9 procesos de PM4DEV, con líneas base de alcance, tiempo y costo, gestión de riesgos y control con valor ganado. & Clyde N. Paricahua Pari \\
Qué falta para el MVP final? & Cuentas en Supabase, punto de referencia en el mapa, criterio de baño, perfil y el piloto con 10 arrendatarios y 5 propietarios. & Clyde N. Paricahua Pari \\
\bottomrule
\end{tabularx}
}

\Needspace{14\baselineskip}
\section{Lista de Verificación}
{\small
\rowcolors{2}{white}{GrisFila}
\noindent\begin{tabularx}{\textwidth}{@{}C{0.8cm} L P{3.0cm} C{1.8cm}@{}}
\toprule
\rowcolor{AzulClaro}
\textbf{} & \textbf{Tarea} & \textbf{Responsable} & \textbf{Cuándo} \\
\midrule
$\square$ & Diapositivas terminadas y revisadas contra los documentos & Viviana M. Quisbert Quispe & Día anterior \\
$\square$ & Ensayo completo con cronómetro (máximo 15 minutos) & Todo el equipo & Día anterior \\
$\square$ & Celular con la app instalada, cargado y con los 5 cuartos de prueba & Eliseo Tarqui Ajahuana & Día anterior \\
$\square$ & Video de respaldo de la demostración & Eliseo Tarqui Ajahuana & Día anterior \\
$\square$ & Expediente impreso o en PDF disponible para consultas & Ruth K. Apaza Solis & Día de la exposición \\
$\square$ & Llegar 15 minutos antes y probar el proyector & Clyde N. Paricahua Pari & Día de la exposición \\
\bottomrule
\end{tabularx}
}

\section{Reglas para la Exposición}
\begin{itemize}[leftmargin=*]
    \item Cada integrante expone su parte y responde las preguntas de su área; el Project Manager modera.
    \item Se usan cifras de los documentos aprobados; si una pregunta no tiene respuesta inmediata, se anota y se responde por escrito.
    \item Las diapositivas llevan pocas palabras y una idea por lámina; el detalle está en el expediente.
\end{itemize}

\Needspace{24\baselineskip}
\section{Firmas}
\noindent En señal de conformidad con el presente documento, los integrantes del equipo firman a continuación:

\vspace{1.5cm}
\begin{center}
\begin{tabular}{c c c}
    \parbox{6.5cm}{\centering \rule{6.0cm}{0.6pt}\\[5pt] \textbf{Clyde Neil Paricahua Pari}\\[2pt] {\footnotesize Project Manager}}
    & \hspace{0.8cm} &
    \parbox{6.5cm}{\centering \rule{6.0cm}{0.6pt}\\[5pt] \textbf{Ruth Karina Apaza Solis}\\[2pt] {\footnotesize Analista de Planificación}}
    \\[2.0cm]
    \parbox{6.5cm}{\centering \rule{6.0cm}{0.6pt}\\[5pt] \textbf{Eliseo Tarqui Ajahuana}\\[2pt] {\footnotesize Especialista Técnico / Arquitectura}}
    & \hspace{0.8cm} &
    \parbox{6.5cm}{\centering \rule{6.0cm}{0.6pt}\\[5pt] \textbf{Viviana Mary Quisbert Quispe}\\[2pt] {\footnotesize Gestora de Comunicación, Riesgos y Documentación}}
\end{tabular}
\end{center}
\end{document}
